\section{Thesis Structure}\label{sec:thesis-structure}

This dissertation is organised as follows. \emph{Chapter~\ref{ch:introduction}}
(the present chapter) establishes the problem, the research gap, the aim,
objectives, research questions and scope of the study.
\emph{Chapter~\ref{ch:literature-review}} reviews related work on LLM copilots
for security, adversarial attacks on LLM-integrated applications, and
hallucination-detection and uncertainty methods, and situates the guarded
intent-recognition approach within it.
\emph{Chapter~\ref{ch:methodology}} presents the research approach: the
objectives-to-method mapping, the design of the evaluation data set, the
metrics, and the experiment and ablation protocol used to answer RQ1--RQ3.
\emph{Chapter~\ref{ch:product-design}} develops the product and solution
design of the Intent Recognition Engine, including the pipeline stages, the
security countermeasures and the intent-contract interface.
\emph{Chapter~\ref{ch:implementation-testing}} describes the implementation,
the integration of the two offline backends, and the automated test suite.
\emph{Chapter~\ref{ch:results}} reports the results of the pipeline-value,
layer-ablation and backend-comparison experiments and interprets them against
the research questions.
\emph{Chapter~\ref{ch:discussion}} discusses the findings, the threats to
validity, and the residual failure modes of the guarded design.
\emph{Chapter~\ref{ch:commercialisation}} considers the commercialisation and
sustainability of the approach.
\emph{Chapter~\ref{ch:conclusion}} concludes the dissertation, restates the
contributions, and outlines directions for future work. Supplementary material
appears in the appendices: the data-set and evaluation instruments
(Appendix~\ref{app:a}) and the declaration on academic integrity and the use of
artificial-intelligence tools (Appendix~\ref{app:b}).